\chapter{The Programme, the History and the Results}\label{ch:programme}\index{open problems}

This last chapter sums up the book in three parts: a research programme in twelve numbered problems (with one added, 1a), ordered by our estimate of difficulty; the history of the ideas, with the record of the existence result announced in 2026; and the list of the results of the book with their status, followed by a scorecard of the ladder against the data and the conjectures that remain. The status of each result is the one stated in its home chapter, which is cited in parentheses.

\section{A Research Programme in Twelve Problems}\label{sec:programme}

The twelve problems below come with their statement, the evidence, the tools, and a first step. Problems 1 to 4 are the most concrete: 1 (for the projected ladder) and 3 are solved and 4 is answered numerically. The last is the Clay problem.

\subsection{Certifiable now}

\paragraph{\normalfont\scshape 1. Certify the ladder on hierarchical lattices.} Solved for the projected ladder: Corollary~\ref{cor:contracts} proves $R(\beta)<\beta$ for SU(2), for every $\beta$ and every scale factor, from Hamilton's Harnack inequality~\cite{Ha93}, the maximum principle and an explicit estimate in the initial layer (Proposition~\ref{prop:initial}), with no computer assistance. What remains is the same statement for SU(3), and for the exact ladder without projection.

\paragraph{\normalfont\scshape 1a. The exact ladder on counts.} By Theorem~\ref{thm:coalescence} the exact ladder is a process on integer counts. Its coalescence kernel is hypergeometric and totally positive in the count and either parent (Corollary~\ref{cor:hypergeometric}); prove that log concavity of the count is preserved (Open problem~\ref{op:logconcave}). By Theorem~\ref{thm:szegocriterion} it would follow from the smoothing statement of Conjecture~\ref{conj:smoothing}, which concerns one law at a time. The stronger and cleaner target is Conjecture~\ref{conj:Lclosed}: the exact ladder preserves the class of link densities with only real zeros in $\cos\theta$, all at most $-1$. Heat flow is proved to preserve it (Theorem~\ref{thm:zeros}); what is missing is a theorem of Malo type~\cite{Malo} for zeros of even multiplicity (Open problem~\ref{op:malo}). Monotonicity is a separate matter: the ladder is not monotone from every start (Table~\ref{tab:notmonotone}).

\paragraph{\normalfont\scshape 2. The string tension function on hierarchical lattices.} Existence and positivity of $\sigma a^2(\beta)=\lim b^{-2m}(-\log u_2^{(m)})$, given that the flow reaches strong coupling, is now Proposition~\ref{prop:sigmaexists}. For the projected ladder the hypothesis holds at every coupling by Corollary~\ref{cor:contracts}; what remains is monotonicity in $\beta$, and the hypothesis for the exact ladder, which is the second half of problem 1.

\paragraph{\normalfont\scshape 3. The second order coefficient of the ladder.} Solved: Theorem~\ref{thm:expansion} gives the expansion in closed form, $\frac{b^2-1}{3b^2}+\frac{b^2-1}{9b^4\beta}+\frac{(b^2-1)(2b^2-1)}{54b^6\beta^2}$, and the scheme independent ratio $1/(6b^2)$ against the true $17/(22\pi^2)$ for SU(2). Theorem~\ref{thm:expansionN} gives the first two orders for every SU(N), with ratio $N/(12b^2)$.

\paragraph{\normalfont\scshape 4. The Wilson law under the ladder.} Answered numerically: by the local heat kernel time of Proposition~\ref{prop:localtime}, the ladder from the Wilson law at $\beta=2.5$ is within $1.2$ percent of the heat kernel family in the likelihood ratio order after two rungs. Theorem~\ref{thm:tp2} gives the condition under which decimation preserves that order, monotonicity of $\Delta f/f$; heat kernels satisfy it and bond moved laws fail it only at the antipode. Theorem~\ref{thm:hemisphere} and Proposition~\ref{prop:truncated} give monotonicity on the near hemisphere with an exponentially small error at weak coupling, and, numerically and not yet proved, at strong coupling ($\beta\le0.69$); the window $0.7<\beta<2$ remains.

\subsection{Hard but concrete}

\paragraph{\normalfont\scshape 5. The likelihood ratio property for SU(3).} The elementary bound of Proposition~\ref{thm:mlrlarge} is proved for SU(2) at large time ($t\ge3.2140$); we expect it to extend to any compact group, but a proof is open. At small time the heat kernel on the maximal torus is a Gaussian sum over the coroot lattice divided by the Weyl denominator. Is there a product formula in two variables? First step: test numerically whether the heat kernel on SU(3) is multivariate totally positive of order two in the appropriate order on the alcove.

\paragraph{\normalfont\scshape 6. A Griffiths inequality for SU(2).} Solved for expectations: Theorem~\ref{thm:doubling} gives $W_{12}\ge W_{11}^2$ and $W_{22}\ge W_{11}^4$ for every compact group, coupling, dimension and character, from reflection positivity. What remains is the same comparison in the likelihood ratio order.

\paragraph{\normalfont\scshape 7. Domination of the Migdal Kadanoff law.} Now a sharp conjecture: $W_{22}(\beta)\le u(2^{d-2}\beta)^4$ with $u=I_2/I_1$. Theorem~\ref{thm:heatbath} proves it for SU(2) with $2(d-1)$ in place of $2^{d-2}$. It holds in ten simulations, at strong coupling, and at weak coupling to leading order, with room in four dimensions and by four percent in three. First step: improve the heat bath bound by integrating two links of a plaquette jointly, which removes the double counting of their common staples.

\paragraph{\normalfont\scshape 8. The origin of $11/3$ in the ladder.} Explain why the ladder at $b=2$ reproduces $97$ percent of the one loop coefficient from the curvature term alone, and find the modification that restores the missing three percent and the missing half of the two loop coefficient together.

\subsection{Research problems}

\paragraph{\normalfont\scshape 9. A uniform log Sobolev inequality at one rung.} If Conjecture~\ref{prop:slicecurv} holds, it gives constant $2$ for a single sliced link. Prove a log Sobolev inequality for the whole sliced conditional, all links given all slices, with a constant that depends on the slices only through the geometry of their caps. This would turn the dynamic route of Section~\ref{ch:dynamic} into a statement about the law of the slice levels alone.

\paragraph{\normalfont\scshape 10. The mass gap in three dimensions.} The problem of Section~\ref{ch:threed}: $m(\beta)\ge c/\beta$ for three dimensional SU(2) at every coupling. Evidence: simulations, the Karabali, Kim and Nair picture~\cite{KKN}, the three dimensional ladder low by twenty six percent, and the sandwich at eight couplings. The flow is contracting, so a fixed multiplicative error per step is affordable.

\paragraph{\normalfont\scshape 11. The comparison in four dimensions for an adapted action.} Choose the lattice action and the block spin rule to make the comparison as easy as possible: heat kernel plaquette weights, which are closed under decimation, and a block rule under which the effective action stays close to a single plaquette law. Prove the sandwich for that pair. Universality then carries the gap to the continuum limit, which does not depend on the lattice action.

\subsection{The Clay problem}

\paragraph{\normalfont\scshape 12. The mass gap in four dimensions.} The Thorin form of the conjecture, Section~\ref{ch:isthere}: a uniform decay rate for the Route B Thorin densities of every gauge invariant observable along the renormalized trajectory. Problems 1, 6, 7 and 11 are steps 3 and 4 of the path; the existence result announced in 2026 and endorsed by a review group, and Ba\l aban's analysis~\cite{Ba87}, bear on step 5.

\subsection{A suggested order of work}

For a group with one analyst and one computational scientist we would begin with problem 1, which is finite, and with the crossover window that problem 4 has become; then problem 7, the Migdal Kadanoff end of the sandwich, now that the independent end is a theorem; then problem 10. Problems 5 and 8 are suited to a student. Problem 12 should be attempted only through the others.

\section{Historical Notes}\label{ch:history}\index{Wilson, K. G.}\index{Creutz, M.}\index{Ba\l aban, T.}\index{Thorin, O.}

The ideas in this book come from three communities that did not often read one another: mathematical physicists working on constructive field theory, probabilists working on infinitely divisible laws and Brownian functionals, and statisticians working on Monte Carlo. These notes trace each line and mark where they crossed.

\subsection{Gauge theory, 1918 to 1974}

Weyl introduced the gauge principle in 1918, as a scale symmetry of the metric, and recast it in 1929~\cite{We29} as the phase symmetry of the electron wave function. Yang and Mills~\cite{YM54} extended it in 1954 to the nonabelian isospin group. Before their paper appeared in October 1954, Pauli had raised an objection at Yang's seminar at the Institute for Advanced Study in February of that year: the quanta of the new field must be massless, and no such particles are seen. The objection is the mass gap problem in its first form.

For a decade the theory was a curiosity. Faddeev and Popov~\cite{FP67} in 1967 showed how to quantise it; 't Hooft in 1971 and 't Hooft and Veltman in 1972~\cite{tH71a,tH71,tHV72} proved it renormalizable. Gross and Wilczek~\cite{GW73}, and Politzer~\cite{Po73}, found asymptotic freedom in 1973: the coupling weakens at short distances, the opposite of electrodynamics. That made the theory a candidate for the strong force, and it turned Pauli's objection inside out. If the coupling grows at long distances, the massless quanta need not appear: they are confined, and what appears are massive bound states.

Wilson's lattice formulation of 1974~\cite{Wi74} made this computable. He replaced the continuum field by group elements on the edges of a lattice, with a weight that is exactly gauge invariant, and showed that at strong coupling the theory confines, by an expansion that statisticians would recognise as a high temperature expansion of a Gibbs measure. Creutz's simulations of 1980~\cite{Cr80} showed, for SU(2) in four dimensions, a smooth passage from Wilson's strong coupling regime to the asymptotically free one. Every attempt on the mass gap since has been an attempt to prove what Creutz saw.

\subsection{Constructive field theory, 1950 to 2026}

Wightman's axioms of the 1950s, and the Euclidean axioms of Osterwalder and Schrader~\cite{OS} in 1973 and 1975, say what it means for a quantum field theory to exist. Reflection positivity, the Euclidean axiom that produces a Hilbert space and a positive Hamiltonian, is the property on which this book rests. Glimm and Jaffe~\cite{GJ}, Nelson~\cite{Ne73}, Simon and others constructed scalar theories in two and three dimensions in the 1960s and 1970s. For gauge theories, Osterwalder and Seiler~\cite{OSe} proved reflection positivity of Wilson's lattice measure and the mass gap at strong coupling in 1978; L\"uscher constructed the transfer matrix in 1977~\cite{Lu77}.

In two dimensions the theory is exactly solvable, by Migdal in 1975~\cite{Mi75} and in the continuum by Gross, King and Sengupta~\cite{GKS89}, Driver~\cite{Dr89}, and Witten~\cite{Wi91} around 1990, as a theory of Brownian motion on the group. In four, Ba\l aban carried out a renormalization group analysis of the lattice theory in a long series of papers from 1983 to 1989~\cite{Ba87}, proving ultraviolet stability. Magnen, Rivasseau and S\'en\'eor gave a different construction in a finite volume in 1993~\cite{MRS}. Neither reached the infrared. The Clay Mathematics Institute made existence and the mass gap one of its seven Millennium Problems in 2000, with a problem description by Jaffe and Witten~\cite{JW}.

On hierarchical lattices, where the Migdal Kadanoff recursion of 1975 and 1976~\cite{Mi75,Ka75,Ka76} is exact, Ito proved confinement for SU(2) in 1985~\cite{Ito85}. Tomboulis proposed in 2007~\cite{To07} to bracket the true theory between Migdal Kadanoff flows; Ito and Seiler~\cite{IS} found a gap in the argument. The probabilistic community entered in the 2010s: Chatterjee's work on large $N$ lattice gauge theory~\cite{Ch19so} and his survey for probabilists~\cite{Ch19}, the stochastic quantisation of Chandra, Chevyrev, Hairer and Shen~\cite{CCHS}, and the strong coupling results of Shen, Zhu and Zhu~\cite{SZZ} by log Sobolev inequalities. In 2026 a proof of existence in four dimensions, produced with a language model and formalised, was announced and endorsed by a review group; we have not checked it, and the mass gap was not part of it.

\subsection{The existence result of October 2026}\label{sec:existence2026}\index{existence of Yang Mills theory}\index{Gross, D. J.}

The other half of the Clay problem was reported to have changed status while this book was being written, and since everything here concerns the half that did not, the record should be set down with its limits.

On 26 October 2026 OpenAI announced that its models had produced a mathematically rigorous existence proof for certain interacting quantum field theories with gauge fields and fermions in four dimensions, including quantum chromodynamics and pure Yang Mills theory, building on Ba\l aban's approach \cite{OAI}. By its account the natural language proof of Yang Mills existence dates from 9 September and a formalisation in Lean from 16 October. A working group chaired by David Gross, with Michael Douglas, Arthur Jaffe, Igor Klebanov and Nikita Nekrasov, took charge of it. On 9 November the group stated that the submission proves the existence of pure Yang Mills theory under the standard statements of the problem, and on 10 November it released the Lean code together with its own validation document. The written proof, of more than a thousand pages, had not been released when this was written. The announcement does not claim confinement or the mass gap, which remain open.

Three cautions. We have read the announcements and reports of them, not the proof: the existence result is announced and endorsed by a review group, and it has not been independently checked, by us or by refereed publication. The assessment is that of a group chosen for the purpose, the review has been disputed in public, among others by Woit \cite{Woit26}, and refereed publication lies ahead. And what exactly is constructed matters for what follows: Woit argues that existence is obtained by compactness, as a limit along subsequences, which is weaker than what the Jaffe Witten statement of the problem asks for and leaves uniqueness open \cite{Woit26}. Woit's posts also discuss a separate programme of Alp\"oge and Sussman, who posted long manuscripts on three dimensions and a preliminary outline for four~\cite{Woit26}.

What it means for this book is the following.

\begin{enumerate}
\item The division of the problem into an ultraviolet half and an infrared half, used since Chapter~\ref{ch:problem}, is now the division between what is claimed and what is open. Everything in this book, the ladder, the latent count, the hitting times, is about the infrared half.
\item A construction in the manner of Ba\l aban proceeds by block averaging and produces effective actions on coarser lattices. That is the setting of our ladder with the true blocking map in place of the Migdal Kadanoff one. The comparison that is missing in Chapter~\ref{ch:sandwich} can therefore be asked of an object that, if the claim stands, exists: is the effective plaquette law after $n$ blocking steps below the ladder's in the likelihood ratio order?
\item The count of Proposition~\ref{prop:observable}, $E[1/(1-\frac12\operatorname{tr}U_p)]-2$, is defined for any plaquette law and so for each effective action. On the ladder it falls by a known amount per rung. Its behaviour along a rigorous sequence of effective actions would be a quantitative test of whether the ladder describes the true flow; Section~\ref{sec:countlattice} tries it on simulated lattices.
\item Uniqueness of the vacuum and the gap are one question seen at two strengths. A gap gives exponential clustering and with it a unique vacuum; a construction that leaves uniqueness open has left the gap open at its first step. In the language of Chapter~\ref{ch:problem}, the reflection positive measure exists, if the claim stands, and the question is the spectrum of its transfer operator near zero.
\end{enumerate}

Nothing in this book depends on the existence result, and nothing in it is confirmed or refuted by the result. If the proof is sustained, the problem that remains is the one this book is about.

\subsection{Infinitely divisible laws, 1929 to 2026}

De Finetti, Kolmogorov, L\'evy and Khintchine characterised infinitely divisible laws between 1929 and 1937. Bernstein's theorem on completely monotone functions~\cite{Be29} dates from 1929. Thorin introduced generalized gamma convolutions in 1977, to prove that the Pareto law~\cite{Th77} is infinitely divisible, and the lognormal law~\cite{Th77b} followed the same year; Bondesson~\cite{Bo92} developed the theory in the 1980s, with hyperbolic complete monotonicity as its sharpest tool. Biane, Pitman and Yor~\cite{BPY} in 2001 collected the laws that tie the Jacobi theta function and the Riemann zeta function to Brownian excursions; their variables $S$ and $C$ appear in this book as curvature corrections.

Statisticians met the same laws through data augmentation. Scale mixtures of normals go back to Andrews and Mallows~\cite{AM74} in 1974; the horseshoe prior and P\'olya Gamma augmentation~\cite{PSW13} are recent members of the family. The $2\sin^2$ lemma came out of the attempt to apply this toolkit to the reciprocal $\Xi$ function.

\subsection{Monte Carlo, 1949 to 2026}

Metropolis and Ulam named the method in 1949; the Metropolis algorithm for hard discs~\cite{MR53} appeared in 1953. Creutz's heat bath for SU(2) dates from 1980, and Kennedy and Pendleton's version from 1985~\cite{KP85}. Swendsen and Wang introduced cluster updates in 1987~\cite{SW87}, Edwards and Sokal explained them in 1988~\cite{ES88}, and Besag and Green~\cite{BG93}, and Higdon~\cite{Hi98}, brought them into statistics in the 1990s. Besag had given the theory of Markov random fields and pseudolikelihood in 1974 and 1975~\cite{Be74,Be75}, for agricultural field trials. Jerrum and Sinclair's polynomial time scheme for the ferromagnetic Ising model is from 1993~\cite{JS93}. Skilling's nested sampling is from 2004 and 2006~\cite{Sk06}, and the vertical likelihood unification from 2014~\cite{PSvl}.

The two Monte Carlo traditions, of lattice field theory and of Bayesian statistics, solved many of the same problems twice under different names; Section~\ref{ch:dictionary} gives a dictionary.

\subsection{Where the lines cross}

Kac crossed them first, and more often than anyone. In 1951 he recognised P\'olya's lemma on the $\xi$ function in the conjecture of Lee and Yang and proved its first case, as he later recorded in his commentary~\cite{Ka74} in P\'olya's collected papers (Chapter~\ref{ch:kacyang}). His 1949 formula made the Schr\"odinger semigroup an expectation over Brownian paths; his 1966 lecture asked what a spectrum determines; his essay on models~\cite{Ka69} taught that an exactly solvable caricature can carry the truth of a hard problem. Symanzik and Nelson made Euclidean field theory a part of probability. Newman connected the Lee Yang theorem~\cite{LY52} to the Riemann Hypothesis in 1976~\cite{Ne76}. The present book adds one more crossing: the Thorin measure, invented for the Pareto and lognormal laws, is the spectral measure of a quantum field after an offset, and the mass gap is its tail.

\section{Summary of New Results}\label{ch:results}

This section lists the results of the book that we believe are new, with their status and location, followed by the computational findings and the conjectures.

\subsection{Theorems}

\begin{enumerate}
\item The spectral measure of a gauge invariant observable, after a Euclidean offset, is the Thorin measure of a generalized gamma convolution, and the mass gap is its exponential tail index (Chapter~\ref{ch:routeA}).
\item The Polson $2\sin^2$ lemma transports the spectral measure to a nonnegative Thorin density that equals the L\'evy exponent of a symmetric compound Poisson process with the energies as jumps, divided by $\pi t$ with $t$ the Euclidean time, and equals the real time return deficit of a Euclidean smeared state (Section~\ref{ch:routeB}).
\item The gap is the exact exponential decay rate in Euclidean time of these Thorin densities, and their local rates are monotone upper bounds (Section~\ref{ch:routeB}).
\item A uniform Thorin gap inequality with some constant is equivalent to the same inequality with constant one, to a uniform spectral gap, and to uniform exponential decay of correlations (Section~\ref{ch:routeB}).
\item A correlation whose spectrum has two isolated levels is not a generalized gamma convolution in Euclidean time, hence not hyperbolically completely monotone (Section~\ref{ch:routeB}).
\item In two dimensional SU(2) lattice gauge theory, $\Delta(a)\ge\Delta_{\mathrm{cont}}\operatorname{arcsinh}(x)/x$ with $x=3g^2a^2/8$, so the uniform Thorin gap inequality holds from the coarsest lattice to the continuum (Chapter~\ref{ch:twod}).
\item Riemann's kernel $\Phi$ is a combination, with coefficients depending on $u$, of the return density of Brownian motion on SU(2) and its time derivative, equivalently of the sphere partition function of two dimensional SU(2) Yang Mills and its area derivative (Section~\ref{ch:riemann}).
\item The Migdal Kadanoff recursion preserves positive definiteness of plaquette laws at every rung (Section~\ref{ch:mk}).
\item The heat kernel projected Migdal Kadanoff map is monotone whenever the SU(2) heat kernel has monotone likelihood ratio in time (Section~\ref{ch:verification}).
\item Product formula: as a function of $\cos\theta$ the SU(2) heat kernel is $K_t(\pi/2)\prod_m(1+\cos\theta/\gamma_m(t))$ with every $\gamma_m>1$ and increasing in $t$. Hence the heat kernel has strictly monotone likelihood ratio in time at every time, is log concave in $\cos\theta$, and the projected ladder map is monotone at every coupling (Chapter~\ref{ch:progress}).
\item Weak coupling expansion of the ladder in closed form: $\beta-R(\beta)=\frac{b^2-1}{3b^2}+\frac{b^2-1}{9b^4\beta}+\frac{(b^2-1)(2b^2-1)}{54b^6\beta^2}+O(\beta^{-3})$; its scheme independent two loop ratio is $1/(6b^2)$ against the true $17/(22\pi^2)$ for SU(2), so no scale factor matches both loop coefficients (Chapter~\ref{ch:progress}).
\item Total positivity of convolution: if $\Delta f/f$ is nondecreasing in the angle, with $\Delta$ the radial Laplacian of the three sphere, then convolution with the class function law $f$ is totally positive of order two and preserves the likelihood ratio order. Heat kernels satisfy the condition, Wilson laws exactly when $\kappa\le3/2$, and bond moved laws fail it at the antipode, so log concavity in $\cos\theta$ is not sufficient (Chapter~\ref{ch:progress}).
\item Ground state representation: the lattice mass gap is $g^2/(2a)$ times the reciprocal Poincar\'e constant of the ground state measure on the product of groups, the relaxation rate of Nelson's ground state diffusion (Section~\ref{ch:h6ym}).
\item For every SU(N) the ladder's weak coupling expansion is one map in the variable $Nt/12$, to two orders; its two loop ratio is $N/(12b^2)$, and its ratios to the true one and two loop coefficients, $0.971$ and $0.532$ at $b=2$, do not depend on $N$ (Chapter~\ref{ch:progress}).
\item Bond moved heat kernel laws satisfy the total positivity condition on the near hemisphere at every coupling; one step of the ladder is monotone in the likelihood ratio order up to an error bounded by the mass outside a ball (Theorem~\ref{thm:hemisphere}, Proposition~\ref{prop:truncated}).
\item Doubling inequality: for every compact gauge group, coupling, dimension and irreducible representation, $(\langle\chi_\rho(U_C)\rangle/d_\rho)^2\le\langle\chi_\rho(U_{C\#rC})\rangle/d_\rho$ for a loop $C$ in a half space that starts and ends on the reflecting plane, by reflection positivity. Hence $W_{12}\ge W_{11}^2$ and $W_{22}\ge W_{11}^4$: the independent end of the sandwich holds for every character. The statement is of a classical type, in the family of reflection positivity loop inequalities, and what is ours is its form for every character and representation; we had wrongly listed it as open (Section~\ref{ch:griffiths}).
\item Heat bath bound, for SU(2) (a classical multihit estimate, Parisi, Petronzio and Rapuano; what is ours is the explicit bound below and its use in the sandwich): $|W_C|\le\langle\prod_{\ell\in A}u(\beta|V_\ell|)\rangle\le u(2(d-1)\beta)^{|A|}$ for links of the loop sharing no plaquette. With the doubling inequality this is a rigorous sandwich for expectations, $W_{11}^4\le W_{22}\le u(2(d-1)\beta)^4$, whose two ends are caricatures of Migdal Kadanoff type with bond moving powers $1$ and $2(d-1)$; the ladder's power $2^{d-2}$ lies between (Section~\ref{ch:griffiths}).
\item Kadanoff's bound evaluated: the ladder's pressure is a rigorous upper bound on the true pressure at every coupling, numerically $1.63$ against $0.78$ at $\beta=2.5$ in four dimensions, with a gap that appears to grow like $\log\beta$. The naive extension of the one step loop bound to $4\times4$ loops fails numerically, as perimeter behaviour requires (Section~\ref{ch:griffiths}).
\item One plaquette as a discrete system: the reciprocal of the SU(2) plaquette partition function, $\sqrt{2\lambda}/(2I_1(\sqrt{2\lambda}))$, is the Laplace transform of the hitting time of the group $S^3$ by Brownian motion in $\R^4$ (the Laplace transform of this hitting time is classical, Ciesielski and Taylor 1962 and Kent 1978; the identification with the SU(2) plaquette is ours), a generalized gamma convolution with Thorin measure at half the squared zeros of $J_1$; the reciprocal heat kernel in $\cos\theta$ is one with atoms $\gamma_m(t)$ (Chapter~\ref{ch:kacyang}).
\item Zeros on fixed rays: for two dimensional SU(2) with the heat kernel action on a torus of $A$ plaquettes, the Fisher zeros lie at $t=4(s_j+2\pi ik)/A$ with $s_j=-\log q_j$ for the zeros $q_j$ of $1+q^3+q^8+\cdots$, hence on the same rays for every volume, and there are no zeros with $|q|\le0.83$. The sector angle is a numerical result, listed below. The method, the smallest discrete system in one variable, is Kac's (Chapter~\ref{ch:kacyang}).
\item Local heat kernel time: every class function law with negative score is bracketed in the likelihood ratio order by the heat kernels at the infimum and supremum of its local time; convolution with a heat kernel preserves the order (Chapter~\ref{ch:progress}).
\item On the hierarchical lattice the string tension $\lim b^{-2m}(-\log u_2^{(m)})$ exists and is positive whenever $u_2^{(m)}\to0$, that is, once the flow reaches strong coupling (Chapter~\ref{ch:progress}).
\item The heat flow of the Wilson law $e^{\kappa\cos\theta}$ has monotone likelihood ratio in time if and only if $\kappa\le3/2$ (Chapter~\ref{ch:progress}).
\item Slicing each plaquette factor of an SU(2) link gives, when every slice level $c_p$ is positive, a conditional that is uniform on a geodesically convex set ; that it satisfies a log Sobolev inequality with constant $2$, independent of the coupling, is Conjecture~\ref{prop:slicecurv}, open because the Neumann form of the Bakry \'Emery criterion and the smoothing step are missing (Section~\ref{ch:slice}).
\item The ladder contracts: $t\,\partial_x\log K_t$ with $x=\cos\theta$ is strictly decreasing in $t$ for every $t>0$, so the bond moved law $K_t^k$ lies strictly below $K_{t/k}$ in the likelihood ratio order and $R(\beta)<\beta$ for every $\beta$ and every scale factor. The proof uses Hamilton's matrix Harnack inequality, the maximum principle and an explicit small time estimate by images; it is not computer assisted (Theorem~\ref{thm:scaledmlr}, Proposition~\ref{prop:initial}, Corollary~\ref{cor:contracts}).
\item The ladder in continuous scale is a Kushner Stratonovich equation without noise, and its projected beta function is a covariance, $dt/d\varsigma=(4-d)t+(d-2)\mathcal B(t)$ with $\mathcal B>0$; and, if the remainder in Theorem~\ref{thm:expansion} is $O((k-1)t^4)$, $\mathcal B=\frac{t^2}6+\frac{t^3}{36}+\frac{t^4}{432}+\cdots$ (Proposition~\ref{prop:contladder}).
\item A latent count for the heat kernel: $K_t/K_t(\pi/2)=\E[(1+\cos\theta)^{M_0}]$ with $M_0$ a sum of independent Bernoulli variables of parameters $1/\gamma_m$; given the count, $\cos^2(\theta/2)$ is Beta; the Wilson law is the Poisson case (Proposition~\ref{prop:latentcount}).
\item The passage time of the walk on representations of SU(2) to dimension $N$ is a generalized gamma convolution with Thorin atoms $1-\cos(j\pi/N)$ and mean $(N^2-1)/3$ (Proposition~\ref{prop:reppassage}).
\item Coalescence: on SU(2), $f_a*f_b=\sum_jQ(j\mid a,b)f_j$ for $f_j\propto(1+\cos\theta)^j$, with $Q>0$ and mean $ab/(a+b+2)$; the family of mixtures of the $f_j$ contains the heat kernels and the Wilson laws and is closed under both operations of the ladder, so the exact Migdal Kadanoff ladder is a superposition and coalescence process on integer counts and every rung is an absolutely monotone function of $1+\cos\theta$ (Theorem~\ref{thm:coalescence}, Corollary~\ref{cor:Mclosed}).
\item Coalescent duality: in $y=\cos^2(\theta/2)$ Brownian motion on SU(2) is a Wright Fisher diffusion with mutation parameter $\theta_{\mathrm m}=3$; heat flow acts on the $f_j$ by the death process of Kingman's coalescent~\cite{Ki82} with rates $\frac12m(m+2)$, and the latent count of $K_t$ is its number of lineages at time $t/2$. Consequences: a second proof of the likelihood ratio property from total positivity of the death process, and the weight of Haar measure in $K_t$ as the distribution function of a generalized gamma convolution with Thorin atoms at the Casimir values (Theorem~\ref{thm:kingman}, whose consequences are the text after it; the Wright Fisher identification is Proposition~\ref{prop:wf}).
\item Real zeros: convolution on SU(2) is a Hadamard product of Laurent coefficients, and by Malo's theorem~\cite{Malo} heat flow preserves the class of densities $\prod_m(1+\cos\theta/\gamma_m)$ with $\gamma_m\ge1$; the coalescence $f_a*f_b$ has only real zeros at most $-1$, so its count is a sum of independent Bernoulli variables (Theorem~\ref{thm:zeros}, Corollary~\ref{cor:coalzeros}).
\item The coalescence kernel is hypergeometric: $Q(j\mid a,b)=\binom aj\binom{b+2}{b-j}/\binom{a+b+2}{b}$, the law of the intersection of two random subsets of sizes $a$ and $b$ in a set of size $a+b+2$, by the Chu Vandermonde identity; hence $Q$ is totally positive in the count and either parent, and the exact ladder is an urn scheme (Theorem~\ref{thm:hypergeometric}).
\item The Wilson law is a mixture of heat kernels over a tilted Hartman Watson time; the reciprocal of a link density with real zeros is the Laplace transform of a generalized gamma convolution with integer Thorin measure on $[1,\infty)$, and two Wilson laws convolve to $2I_1(r)/r$ with explicit zeros (Chapter~\ref{ch:zeros}).
\item Hitting times bracket the gap: with $h_*$ the least limiting hazard rate over sets of measure at least one half, $h_*\le\lambda_1\le2h_*$; for Brownian motion on SU(2) the hazard rate of the hitting time of the equator tends to the gap $\frac32$ of Brownian motion with generator $\frac12\Delta_{S^3}$ ($\frac34$ in the time $t$ of $K_t$), with an explicit survival series (Chapter~\ref{ch:hitting}).
\item Matching is maximum likelihood: the total action is sufficient for $\beta$, loop sums are sufficient for loop couplings, and the maximum likelihood estimate matches their expectations. Parisi's effective coupling~\cite{Pa80} and Swendsen's matching~\cite{Sw79} are moment estimators, not likelihood equations of the Wilson model. The Wilson family is conjugate for bond moving and the heat kernel family for decimation, and the two families have no nontrivial member in common (Section~\ref{ch:plgbc}).
\item The Migdal Kadanoff step is a Bayes filter recursion, an update by the $(k-1)$st power of the current density followed by a convolution, so the ladder is a nonlinear Markov chain of McKean type and can be run with particles (Section~\ref{ch:h6}).
\end{enumerate}

\subsection{Leading order derivations}

\begin{enumerate}
\item The Migdal Kadanoff decrement in four dimensions is $(1-b^{-2})N^2\E S_{\mathrm{BPY}}/4$ for SU(N), where $S_{\mathrm{BPY}}$ is the Biane Pitman Yor variable with $\E S_{\mathrm{BPY}}=1/3$; its ratio to the one loop coefficient is $\pi^2(1-b^{-2})/(11\ln b)$, equal to $0.971$ at $b=2$ for every $N$ (Section~\ref{ch:mk}).
\item In general dimension, $1/t''=b^{d-4}/t-(b^{d-2}-1)\E S_{\mathrm{BPY}}/(2b^2)$, so the coupling is relevant below four dimensions and marginal in four, with asymptotic freedom produced entirely by Haar curvature (Section~\ref{ch:mk}).
\item The first passage time $J(\beta)$ of the renormalization walk of the projected ladder is $\beta/\delta_\infty+O(\log\beta)$ given a positive limiting decrement $\delta_\infty$ (Proposition~\ref{prop:criterionB}). If in addition the exact recursion is described by the projected map, the strong coupling theory has a positive mass, and $a(\beta)b^{J(\beta)}$ is bounded, the correlation length is at most $b^{J}$ in lattice units and the asymptotic freedom scaling follows; none of these three conditions is proved, and the third, as stated, fails for the ladder at $b=2$, whose decrement differs from the one loop value (Section~\ref{ch:verification}).
\item The one loop coefficient of the ladder is the first Cram\'er coefficient $\frac16$ of the heat kernel on the group, the curvature term of $\log(\theta/\sin\theta)$, found by matching only the quadratic part of the exponent (Remark~\ref{rem:cramerloop}).
\end{enumerate}

\subsection{Computational findings}

\begin{enumerate}
\item The decrement inequality $R(\beta)<\beta$, now Corollary~\ref{cor:contracts}, was first seen on $[0.3,300]$ by a monotone covering with margin $0.051$, and the likelihood ratio property holds on an $800\times1500$ grid (Section~\ref{ch:verification}).
\item A full SU(3) Migdal Kadanoff computation confirms the decrement $9/16$ (Section~\ref{ch:mk}).
\item Migdal Kadanoff trajectories started at weak and intermediate coupling follow, until near the origin, the heat kernel family $u_3=u_2^{8/3}$ (Section~\ref{ch:rg}).
\item The gap is identified sharply as a support constraint in positive deconvolution even when the spectrum is not (Chapter~\ref{ch:estimation}).
\item Nested sampling recovers the two dimensional partition function at every coupling from one run, with root mean square error $0.22$, while the vertical likelihood sampler computes it with $0.03$, path sampling $0.06$, split sampling $0.14$, and the harmonic mean, whose error changes sign between seeds, is worst at $0.27$ (Sections~\ref{ch:latent} and~\ref{ch:vertical}).
\item Against SU(2) string tension data at $2.2\le\beta\le2.635$, the ladder predicts $d\log a/d\beta=-2.77$ while the data give $-3.18$: the ladder runs about fifteen percent too slowly, and in its weak coupling limit it lies between one loop perturbation theory and the data (Chapter~\ref{ch:progress}).
\item In Parisi's effective coupling, with plaquettes from our own $8^4$ runs, the ladder predicts the running within ten percent and two loop perturbation theory within four percent; the ladder's two loop coefficient is $0.021$ against the true $0.0403$ (Chapter~\ref{ch:progress}).
\item From the Wilson law at $\beta=2.5$ the ladder enters a $1.2$ percent neighbourhood of the heat kernel family, in the likelihood ratio order, after two rungs (Chapter~\ref{ch:progress}).
\item The ladder's absolute string tension in four dimensions is too small by a factor of about fifty, a shift of $1.45$ in $\beta$; in three dimensions $\beta a\sqrt\sigma$ is $0.98$ to $1.00$ (mean $0.99$, with a modulation of about two percent) against $1.341$, low by twenty six percent; in four dimensions $a\sqrt\sigma=0.0088$, $0.0068$, $0.0053$, $0.0039$ and $0.0028$ at $\beta=2.2$, $2.3$, $2.4$, $2.5$ and $2.635$ (Chapter~\ref{ch:progress}).
\item Numerically, bond moved heat kernel laws satisfy the total positivity condition on the whole group for $\beta\le0.69$; the bad cap at the antipode has radius $\sqrt{(k-1)t/(2k)}$ at small $t$ and mass at most $8\times10^{-4}$ (Chapter~\ref{ch:progress}).
\item Numerically, the sector of half angle $82.9^\circ$ around the positive axis is free of Fisher zeros for two dimensional SU(2) on a torus, with $|\arg s|=82.9^\circ$ and $88.2^\circ$ for the two lowest zeros, $s$ being the zero at $t=4s/A$; this is not certified (Chapter~\ref{ch:kacyang}).
\item Empirical sandwich: in four and three dimensional SU(2), the law of the $2\times2$ and $4\times4$ loop holonomies lies between the law for independent plaquettes and the Migdal Kadanoff law, in all twenty four comparisons, and in the likelihood ratio order on at least $95$ percent of the mass in eighteen of them, the exceptions being laws within noise of Haar measure (Chapter~\ref{ch:sandwich}).
\item Lowest order perturbation theory gives $\log W_{22}/\log W_{11}=2.739$ in four dimensions and $3.120$ in three, strictly inside the bracket, and the simulations approach these values (Chapter~\ref{ch:sandwich}).
\item $W_{12}>W_{11}^2$ for four dimensional SU(2) at eleven couplings from $\beta=1.2$ to $10$, by more than fifty standard errors above the crossover and by at least three below it: the smallest nonabelian Griffiths inequality holds numerically (Section~\ref{ch:griffiths}).
\item Pseudolikelihood recovers the four dimensional coupling to two parts in a thousand; the coupling fitted to a blocked field ranges from $0.36\beta$ to $2\beta$ with the blocking rule, with the ladder in between (Chapter~\ref{ch:sandwich}).
\item Particle learning recovers the decrement of the ladder from sampled plaquettes: posterior median $0.264$, ninety percent interval $[0.226,0.304]$, truth $0.254$ (Section~\ref{ch:plgbc}).
\item A GBC generator trained on $4^4$ simulations matches blocked $8^4$ ensembles to Wilson couplings with median error $0.03$; consistency across loops selects the blocking rule, and at the selected rule one blocking step shifts $\beta$ by $0.38\pm0.02$ at $\beta=2.5$ (Section~\ref{ch:plgbc}).
\item Plaquettes are pairwise nearly independent, with correlations below five percent at every coupling, yet $\log W_{22}/\log W_{11}$ is $2.8$ to $3.5$ and not $4$ (Chapter~\ref{ch:sandwich}).
\item A particle version of the ladder with $2\times10^5$ group elements tracks the exact flow of $u_2$ to within $0.014$ over five steps from $\beta=3$ and $0.023$ from $\beta=2$ (Section~\ref{ch:h6}).
\item The abelian control: for compact U(1) with heat kernel laws the ladder's decrement is exponentially small, $2\times10^{-5}$ at $\beta=1$ and $2\times10^{-23}$ at $\beta=4$, against $\frac14$ for SU(2); the ladder separates the two but cannot produce the Coulomb phase (Section~\ref{ch:missing}).
\item Fisher zeros of two dimensional SU(2) on a torus leave the imaginary axis from four plaquettes on, with nearest approach to the real axis about $3$; the Lee and Yang property available to gauge theory is a zero free neighbourhood of the positive axis (Chapter~\ref{ch:kacyang}).
\item In two dimensions maximum likelihood from the action attains the Cram\'er Rao bound at $\beta=2$ and $4$, pseudolikelihood has efficiency $0.91$ to $0.85$, and a one bit summary $0.53$ to $0.67$; a GBC generator fed the sufficient statistic reproduces the exact posterior quantiles to $0.03$ (Section~\ref{ch:plgbc}).
\item For windows after $t=1$ a hazard rate estimator gives the gap of SU(2) as $1.497\pm0.017$ where the correlator from the same paths and window gives $1.49\pm0.06$; a GBC generator for the gap of a two level correlator has median error $0.018$ against $0.045$ for the effective mass (Chapter~\ref{ch:hitting} and Section~\ref{ch:plgbc}).
\item In three hundred random trials and in exact rungs of the ladder up to degree $48$, convolution preserves the class of class functions with only real zeros at most $-1$ in $\cos\theta$; conjecturally the exact ladder has a Lee Yang property at every rung (Conjecture~\ref{conj:Lclosed}).
\item To leading order at large counts the mean count changes by $-\frac{13}{16}+\frac{15}{16}r$ per rung, with $r$ the ratio of variance to mean, and $r'=\frac5{16}+\frac r{16}$; this gives $-\frac{13}{16}$ from point masses, $-\frac12$ from heat kernels and $+\frac18$ from Wilson laws, as computed (Section~\ref{sec:hypergeometric}).
\item The harmonic law: the mean count of a convolution of $f_j$ is $1/H_1-\frac76+H_2/2H_1^2-H_3/3H_1^3$ up to $O(1/J)$, which proves the drift $-\frac{13}{16}$ from point masses and the decimation half of the drift law (Theorem~\ref{thm:harmonic}). The mean count is the observable $E[1/(1-\cos\theta)]-2$ (Proposition~\ref{prop:observable}), and the first rung changes it by $-\frac{13}{16}$ in the limit from point masses and by $+\frac18$ from Wilson laws (Theorem~\ref{thm:firstrung}, which gives $-\frac{13}{16}+\frac{15}{16}r+o(1)$ on the class $\mathcal L$ and evaluates the limit of $r$ only in these two cases); the heat kernel value $-\frac12$ is a numerical observation, because the limit $r\to\frac13$ it needs is unproved. The Hartman Watson law is a generalized gamma convolution with Thorin distribution function $-\arg I_{iy}(r)/\pi$ at $t=\frac12y^2$, which counts the zeros of P\'olya's function $K_{iy}(r)$; hence the Wilson action is a Thorin mixture of heat kernel actions (Theorem~\ref{thm:hwggc}, Corollary~\ref{cor:wilsonthorin}). One plaquette in the Hamiltonian formulation has gap $\frac34\lambda+O(\lambda^{-3})$ at strong coupling, tending to $\sqrt2$ at weak coupling (Proposition~\ref{prop:oneplaquette}), and numerically the levels a link needs grow like the square root of its count (Open problem~\ref{op:registers}); its vacuum has only real zeros in $\cos\theta$, all below $-1$ (Proposition~\ref{prop:vacuumzeros}, checked to $160$ digits). On an $8^4$ lattice the plaquette has the count of a Wilson law and the $2\times2$ loop that of a heat kernel, and after two blockings on $16^4$ the blocked plaquette has a smaller count than two rungs of the ladder at $\beta=2.3$, $2.5$ and $2.7$, though not after one blocking at $2.7$ (Section~\ref{sec:countlattice}; numerical, small lattices). The convolution $f*g$ of two laws with real zeros has real zeros whenever $f*f_{\deg g}$ or $g*f_{\deg f}$ has an antiderivative with real zeros (Theorem~\ref{thm:szegocriterion}); the hypothesis held in every trial, and its general truth is Conjecture~\ref{conj:smoothing}. Laws with real zeros are products of fermion determinants $\det(1+\alpha_mU)$, their characters are sums of Schur functions of the hopping parameters, and decimation is a sum over matchings of the flavours (Propositions~\ref{prop:determinant} and~\ref{prop:howe}, Theorem~\ref{thm:matching}). The SU(2) convolution is a radial flow of circle convolutions, its flat limit has real zeros, and closure follows if the two convolutions never share a zero along the flow, or if the lower convolutions alternate in sign at the zeros of the circle convolution (Theorems~\ref{thm:radialflow}, \ref{thm:flowclosure} and~\ref{thm:signcriterion}); both held in every case computed. Every law of the class $\mathcal L$ has a log concave count (Proposition~\ref{prop:Llogconcave}).
\item The exact ladder is not monotone on all of $\mathcal M$: from Wilson laws with $\kappa\ge2$ the fundamental coefficient rises at the first rung, and from bimodal laws the mean count rises (Table~\ref{tab:notmonotone}).
\item On the exact ladder in counts the mean count falls by $0.508$ per rung from $\beta=4$, $3$ and $2$, each law of the count is below the previous one in the likelihood ratio order, (Chapter~\ref{ch:counts}); the total positivity of $Q(j\mid a,b)$ checked there for $a,b\le20$ is now Corollary~\ref{cor:hypergeometric}.
\item The Langevin gap of one plaquette grows from $1.5$ at $\beta=0$ to $6.3$ at $\beta=12$ while the Bakry \'Emery bound $1-\beta/2$ is void from $\beta=2$; a one layer attention head inside GBC has quantile error $0.02$ among eight pooled characters, though the grid of attention scales fixes the split between signal and noise by hand (Section~\ref{ch:spde} and Section~\ref{ch:plgbc}).
\item The frame lemma and its corollary: maps that form an approximate tight frame carry subspaces of small codimension to a sum of small codimension, and a form bound on that sum bounds the number of low eigenvalues and the mean of the spectrum (Lemma~\ref{lem:frame}, Corollary~\ref{cor:lowstates}); with sharp constants the codimension is at most $\alpha D$ for a tight frame of isometries (Proposition~\ref{prop:framesharp}). It bounds averages such as a free energy, not a gap.
\item The budget of the ladder: an inequality of type (M) for the count needs one condition, an error per rung below the drift $\frac12$. At weak coupling the count of the four dimensional plaquette is $4\beta$ to leading order, the ladder from it needs about $8\beta$ rungs and asymptotic freedom about $3.9\beta$ blockings, so the inequality has room (Section~\ref{sec:ineqM}).
\item Pseudolikelihood recovers the coupling of two dimensional SU(2) from single configurations with an error of about fifteen percent at $\beta\le2$ and about five percent at $\beta\ge4$ (Section~\ref{ch:slice}).
\item Weighted Monte Carlo on sorted couplings achieves the Yakowitz rate~\cite{Ya78} $n^{-2}$ in thermodynamic integration, a factor of several thousand over plain Monte Carlo at a few hundred rungs (Section~\ref{ch:yakowitz}).
\end{enumerate}

\subsection{The ladder: a scorecard}

\begin{center}
\small
\begin{tabular}{p{2.35in}p{0.8in}p{0.8in}}
\toprule
quantity & ladder & truth \\ \midrule
positivity at every rung & theorem & theorem \\
monotone projected map, every coupling & theorem & expected \\
one loop decrement, SU(2), $b=2$ & 0.2500 & 0.2575 \\
two loop ratio $\rho$ & $1/24$ & $17/(22\pi^2)$ \\
slope of $\log a$ in $\beta$, bare coupling & $-2.77$ & $-3.18$ \\
slope in the effective coupling & $-2.64$ & $-2.40$ \\
$a\sqrt\sigma$ at $\beta=2.5$, four dimensions & 0.0039 & 0.187 \\
$\beta a\sqrt\sigma$, three dimensions & 0.98 to 1.00 & 1.341 \\
$\log W_{22}/\log W_{11}$, weak coupling, $d=4$ & 1 & 2.739 \\
$u_2$ of the $2\times2$ loop at $\beta=2.5$ & 0.665 & 0.261 \\
side of the truth it lies on & stiffer & \\
\bottomrule
\end{tabular}
\end{center}

The ladder is right about structure and first order rates and wrong, always in the same direction, about magnitudes: it makes the blocked theory too stiff. The independent caricature errs in the opposite direction at every entry where it can be computed, though the table does not show it. That consistency of sign is what makes a sandwich plausible.

\subsection{Conjectures and open problems}

\begin{enumerate}
\item The mass gap conjecture in Thorin form (Section~\ref{ch:isthere}).
\item The Thorin sandwich, the dynamic sandwich, and the latent variable coupling as routes to step 4 (Section~\ref{ch:isthere}).
\item The existence of an FPRAS for lattice gauge partition functions at every coupling (Section~\ref{ch:jerrum}).
\item The bare Migdal Kadanoff end of the sandwich, $W_{22}(\beta)\le(I_2(2^{d-2}\beta)/I_1(2^{d-2}\beta))^4$, true in ten simulations, at strong coupling and, to leading order, at weak coupling, and both halves in the likelihood ratio order (Chapter~\ref{ch:sandwich} and Section~\ref{ch:griffiths}).
\item A mass gap for three dimensional SU(2) at all couplings by contracting brackets (Section~\ref{ch:threed}).
\item The origin of the agreement between the ladder's $\E S_{\mathrm{BPY}}=1/3$ and the coefficient $4-\frac13$ (Section~\ref{ch:h6}).
\item The likelihood ratio property for the SU(3) heat kernel at all times; the elementary bound gives it for SU(2) at large time (Chapter~\ref{ch:progress}).
\item An improved ladder whose two loop ratio for SU(2) is $17/(22\pi^2)$ and not $1/(6b^2)$ (Chapter~\ref{ch:progress}).
\item Monotonicity of the ladder step in the window $0.7<\beta<2$, where the bad cap has mass up to $8\times10^{-4}$ (Chapter~\ref{ch:progress}).
\item The class $\mathcal L$ is closed under convolution on SU(2), so that the exact ladder has real zeros at every rung (Conjecture~\ref{conj:Lclosed}); the smoothing statement behind it (Conjecture~\ref{conj:smoothing}); and the weak coupling sandwich (Conjecture~\ref{prop:weaksandwich}, Chapter~\ref{ch:sandwich}).
\item A Poincar\'e inequality for the ground state measure, through Greensite's dimensional reduction~\cite{Gr79} (Section~\ref{ch:h6ym}).
\end{enumerate}

\subsection{Where things stand}

In one paragraph. The mass gap is a statement about the support of a positive measure, and the Thorin transport turns it into a statement about decay rates that can be proved outright in two dimensions and, for the projected ladder, on hierarchical lattices. The proposal is to carry the four dimensional argument by the Migdal Kadanoff ladder; it is a filter recursion on laws on the group, whose heat kernel projection is monotone by a product formula for the heat kernel, and asymptotically free by the curvature of the group. Against lattice data it has roughly the right rate (fifteen percent slow in the bare coupling) and the wrong scale. Simulations place the true block spin law between the ladder and the independent caricature, in the likelihood ratio order on most of the mass, in three and four dimensions, at one and two steps, and lowest order perturbation theory, evaluated numerically, gives that placement at weak coupling (Conjecture~\ref{prop:weaksandwich}). The independent end is now a theorem for expectations, by reflection positivity. A rigorous upper bound of the same form exists for SU(2) with bond moving power $2(d-1)$. What is missing is that bound with the ladder's power $2^{d-2}$, and the likelihood ratio form of both halves.

Nothing in this book proves the mass gap. What it offers is a language in which the remaining step is a comparison of probability laws on a compact group, a set of exact results that show the language is the right one, and a list of problems, ordered by difficulty, each of which is smaller than the one before it was reduced from.
